% =============================================================================
%  Glossario do relatorio
%
%  Entradas para os termos que o texto usa repetidamente e que nao sao obvios
%  para quem chega ao C++ moderno pela primeira vez. Regra de inclusao: o termo
%  aparece em mais de um capitulo E nao e auto-explicativo em portugues.
%
%  No texto, use \gls{chave} para a primeira ocorrencia relevante de um termo
%  numa secao; as demais podem ser \texttt{} normal.
% =============================================================================

% --- O modelo de objeto e de valor -------------------------------------------

\newglossaryentry{raii}{
  sort={RAII},
  name={RAII},
  description={\emph{Resource Acquisition Is Initialization}. O idioma central
    do C++: o tempo de vida de um recurso (memória, arquivo, \emph{lock},
    \emph{socket}) é amarrado ao tempo de vida de um objeto. Adquirir é
    construir; liberar é destruir. Como o destrutor roda também durante o
    desenrolamento da pilha, o idioma é o que torna exceções seguras}}

\newglossaryentry{ub}{
  sort={UB comportamento indefinido},
  name={UB (comportamento indefinido)},
  description={Situação em que o padrão não impõe nenhuma exigência ao
    programa. Não é ``dá erro'': é o compilador autorizado a supor que aquilo
    nunca acontece, e a otimizar a partir dessa suposição. Ler variável não
    inicializada, estourar um índice, violar \emph{strict aliasing} e escrever
    num objeto genuinamente \texttt{const} são casos deste relatório}}

\newglossaryentry{categoriadevalor}{
  sort={categoria de valor},
  name={categoria de valor},
  description={A classificação de uma \emph{expressão} segundo duas perguntas
    independentes --- tem identidade? pode ser movida? --- que produzem
    \emph{lvalue}, \emph{xvalue} e \emph{prvalue}. É do que depende a escolha
    entre copiar e mover, e a fonte da regra de que tudo o que tem nome é
    \emph{lvalue}}}

\newglossaryentry{move}{
  sort={move},
  name={move},
  description={Transferência do conteúdo de um objeto para outro, deixando a
    origem num estado válido porém não especificado. \texttt{std::move} não
    move nada: é um \texttt{static\_cast} para \emph{rvalue reference} que
    apenas \emph{autoriza} a operação, executada pelo construtor ou pela
    atribuição de \emph{move} do tipo}}

\newglossaryentry{elisao}{
  sort={elisao de copia},
  name={elisão de cópia},
  description={Omissão da construção de um temporário: o objeto é construído
    diretamente no destino. Desde o C++17 é \emph{obrigatória} para
    \emph{prvalues} (RVO) e opcional para variáveis locais retornadas (NRVO).
    Escrever \texttt{return std::move(local)} desliga a segunda}}

\newglossaryentry{forwarding}{
  sort={perfect forwarding},
  name={\emph{perfect forwarding}},
  description={Repasse de um argumento preservando tipo, constância e
    categoria de valor. Exige as duas peças juntas: o parâmetro \texttt{T\&\&}
    deduzido (referência universal) e o \texttt{std::forward<T>} na hora de
    repassar}}

\newglossaryentry{slicing}{
  sort={object slicing},
  name={\emph{object slicing}},
  description={Cópia de um objeto derivado para uma variável do tipo base:
    o que não cabe em \texttt{sizeof(Base)} é amputado, e com ele o
    polimorfismo. Nenhum \emph{warning} é emitido. Passar hierarquia por valor
    e usar \texttt{std::vector<Base>} são as duas formas usuais de cometê-lo}}

\newglossaryentry{vtable}{
  sort={vtable vptr},
  name={vtable / vptr},
  description={A implementação usual do despacho dinâmico: cada classe
    polimórfica tem uma tabela de ponteiros de função (\emph{vtable}) e cada
    objeto guarda um ponteiro para a tabela da sua classe real (\emph{vptr}).
    O custo é um ponteiro por objeto e uma indireção por chamada --- mas o
    preço maior costuma ser a impossibilidade de \emph{inline}}}

\newglossaryentry{rtti}{
  sort={RTTI},
  name={RTTI},
  description={\emph{Run-Time Type Information}. O mecanismo que sustenta
    \texttt{typeid} e \texttt{dynamic\_cast}. Exige pelo menos uma função
    virtual na base e pode ser desligado com \texttt{-fno-rtti}}}

% --- Genericidade e compilacao -----------------------------------------------

\newglossaryentry{sfinae}{
  sort={SFINAE},
  name={SFINAE},
  description={\emph{Substitution Failure Is Not An Error}. Se a substituição
    dos parâmetros de \emph{template} produz um tipo inválido na assinatura,
    o candidato é apenas descartado da resolução de sobrecarga em vez de gerar
    erro. Era a forma pré-C++20 de restringir \emph{templates}; \emph{concepts}
    a substitui com mensagens legíveis}}

\newglossaryentry{crtp}{
  sort={CRTP},
  name={CRTP},
  description={\emph{Curiously Recurring Template Pattern}: a classe derivada
    passa a si mesma como argumento de \emph{template} da base
    (\texttt{struct D : Base<D>}). Dá polimorfismo resolvido em compilação ---
    sem \emph{vtable}, sem indireção e com \emph{inline} --- ao preço de
    perder a base comum em tempo de execução}}

\newglossaryentry{concept}{
  sort={concept},
  name={\emph{concept}},
  description={Predicado sobre tipos, avaliado em compilação, que restringe
    quais argumentos um \emph{template} aceita. Move o erro do interior da
    instanciação para o ponto da chamada e participa da resolução de
    sobrecarga por \emph{subsumption}}}

\newglossaryentry{constexpr}{
  sort={constexpr consteval},
  name={\texttt{constexpr} / \texttt{consteval}},
  description={\texttt{constexpr} marca o que \emph{pode} ser avaliado em
    compilação; \texttt{consteval} marca o que \emph{tem} de ser. A diferença
    entre ``o compilador talvez calcule'' e ``o compilador calcula ou recusa
    o programa''}}

\newglossaryentry{typeerasure}{
  sort={type erasure},
  name={\emph{type erasure}},
  description={Guardar objetos de tipos não relacionados atrás de uma
    interface uniforme, apagando o tipo estático mas preservando o
    comportamento. \texttt{std::function}, \texttt{std::any} e o
    \texttt{std::type\_index} do ECS deste repositório são três realizações
    diferentes da mesma ideia}}

% --- Biblioteca padrao -------------------------------------------------------

\newglossaryentry{iterador}{
  sort={iterador},
  name={iterador},
  description={A abstração que separa algoritmos de containers: um objeto que
    aponta para um elemento e sabe avançar. As categorias (\emph{input},
    \emph{forward}, \emph{bidirectional}, \emph{random access}) definem que
    algoritmos podem operar sobre qual container e a que custo}}

\newglossaryentry{view}{
  sort={view},
  name={\emph{view} (C++20)},
  description={Um \emph{range} que não possui os elementos e cujas operações
    de cópia e destruição são O(1). \emph{Views} são preguiçosas: compor
    \texttt{filter} com \texttt{transform} não percorre nada até que alguém
    itere o resultado}}

\newglossaryentry{invalidacao}{
  sort={invalidacao de iteradores},
  name={invalidação de iteradores},
  description={A perda de validade de iteradores, ponteiros e referências
    após uma modificação do container. As regras variam por container e são a
    origem mais comum de \gls{ub} em código que ``funcionava''}}

\newglossaryentry{pmr}{
  sort={pmr},
  name={\texttt{std::pmr}},
  description={\emph{Polymorphic Memory Resources} (C++17). Permite trocar a
    estratégia de alocação de um container em tempo de execução, sem que a
    escolha entre no tipo --- ao contrário do parâmetro \texttt{Allocator}
    clássico, que faz \texttt{vector<int,A>} e \texttt{vector<int,B>} serem
    tipos distintos}}

% --- Concorrencia ------------------------------------------------------------

\newglossaryentry{datarace}{
  sort={data race},
  name={\emph{data race}},
  description={Dois acessos concorrentes ao mesmo endereço, ao menos um deles
    de escrita, sem sincronização entre eles. É \gls{ub} --- não ``valor
    errado ocasional''. O detector prático é o ThreadSanitizer}}

\newglossaryentry{atomico}{
  sort={operacao atomica},
  name={operação atômica},
  description={Operação indivisível sobre uma variável compartilhada, com uma
    ordem de memória associada. Resolve a corrida sobre \emph{aquela}
    variável; não estabelece invariantes entre duas variáveis --- para isso
    ainda é preciso um \emph{mutex}}}

\newglossaryentry{corotina}{
  sort={co-rotina},
  name={co-rotina (C++20)},
  description={Função que pode suspender a própria execução e ser retomada
    depois, preservando o estado local num \emph{frame} alocado à parte.
    A linguagem fornece só o mecanismo (\texttt{co\_await},
    \texttt{co\_yield}, \texttt{co\_return}); os tipos \texttt{generator} e
    \texttt{task} são escritos pelo usuário ou pela biblioteca}}

% --- Ferramentas e ecossistema -----------------------------------------------

\newglossaryentry{meson}{
  sort={Meson},
  name={Meson},
  description={O sistema de \emph{build} deste repositório. Descreve o projeto
    numa DSL declarativa e gera arquivos do Ninja; a compilação em si é do
    Ninja. Todo subprojeto tem um \texttt{meson.build} próprio, registrado em
    \texttt{src/meson.build}}}

\newglossaryentry{conan}{
  sort={Conan},
  name={Conan 2},
  description={O gerenciador de dependências do repositório. Lê o
    \texttt{conanfile.py}, resolve e compila as bibliotecas de terceiros e
    gera os arquivos que fazem o \texttt{dependency()} do Meson encontrá-las}}

\newglossaryentry{abi}{
  sort={ABI},
  name={ABI},
  description={\emph{Application Binary Interface}: o contrato em nível de
    binário --- convenção de chamada, \emph{layout} de objetos, decoração de
    nomes. É o que decide se dois \texttt{.o} compilados separadamente podem
    ser ligados, e o que \texttt{extern "C"} e o ABI System V x86-64 fixam}}
